\section{Relative Proj via glueing}
\label{section-relative-proj-via-glueing}

\begin{situation}
\label{situation-relative-proj}
Here $S$ is a scheme, and $\mathcal{A}$
is a quasi-coherent graded $\mathcal{O}_S$-algebra.
\end{situation}

\noindent
In this section we outline how to construct a morphism
of schemes
$$
\underline{\text{Proj}}_S(\mathcal{A}) \longrightarrow S
$$
by glueing the homogeneous spectra $\text{Proj}(\Gamma(U, \mathcal{A}))$
where $U$ ranges over the affine opens of $S$. We first show that the
homogeneous spectra of the values of $\mathcal{A}$ over affines form a
suitable collection of schemes, as in Lemma \ref{lemma-relative-glueing}.

\begin{lemma}
\label{lemma-proj-inclusion}
In Situation \ref{situation-relative-proj}.
Suppose $U \subset U' \subset S$ are affine opens.
Let $A = \mathcal{A}(U)$ and $A' = \mathcal{A}(U')$.
The map of graded rings $A' \to A$ induces a morphism
$r : \text{Proj}(A) \to \text{Proj}(A')$, and the diagram
$$
\xymatrix{
\text{Proj}(A) \ar[r] \ar[d] &
\text{Proj}(A') \ar[d] \\
U \ar[r] &
U'
}
$$
is cartesian. Moreover there are canonical isomorphisms
$\theta : r^*\mathcal{O}_{\text{Proj}(A')}(n) \to
\mathcal{O}_{\text{Proj}(A)}(n)$ compatible with multiplication maps.
\end{lemma}

\begin{proof}
Let $R = \mathcal{O}_S(U)$ and $R' = \mathcal{O}_S(U')$.
Note that the map $R \otimes_{R'} A' \to A$ is an isomorphism as
$\mathcal{A}$ is quasi-coherent
(see Schemes, Lemma \ref{schemes-lemma-widetilde-pullback} for example).
Hence the lemma follows from
Lemma \ref{lemma-base-change-map-proj}.
\end{proof}

\noindent
In particular the morphism $\text{Proj}(A) \to \text{Proj}(A')$
of the lemma is an open immersion.

\begin{lemma}
\label{lemma-transitive-proj}
In Situation \ref{situation-relative-proj}.
Suppose $U \subset U' \subset U'' \subset S$ are affine opens.
Let $A = \mathcal{A}(U)$, $A' = \mathcal{A}(U')$ and $A'' = \mathcal{A}(U'')$.
The composition of the morphisms
$r : \text{Proj}(A) \to \text{Proj}(A')$, and
$r' : \text{Proj}(A') \to \text{Proj}(A'')$ of
Lemma \ref{lemma-proj-inclusion} gives the
morphism $r'' : \text{Proj}(A) \to \text{Proj}(A'')$
of Lemma \ref{lemma-proj-inclusion}. A similar statement
holds for the isomorphisms $\theta$.
\end{lemma}

\begin{proof}
This follows from Lemma \ref{lemma-morphism-proj-transitive} since
the map $A'' \to A$ is the composition of $A'' \to A'$ and
$A' \to A$.
\end{proof}

\begin{lemma}
\label{lemma-glue-relative-proj}
In Situation \ref{situation-relative-proj}.
There exists a morphism of schemes
$$
\pi : \underline{\text{Proj}}_S(\mathcal{A}) \longrightarrow S
$$
with the following properties:
\begin{enumerate}
\item for every affine open $U \subset S$ there exists an isomorphism
$i_U : \pi^{-1}(U) \to \text{Proj}(A)$ with $A = \mathcal{A}(U)$, and
\item for $U \subset U' \subset S$ affine open the composition
$$
\xymatrix{
\text{Proj}(A) \ar[r]^{i_U^{-1}} &
\pi^{-1}(U) \ar[rr]^{inclusion} & &
\pi^{-1}(U') \ar[r]^{i_{U'}} &
\text{Proj}(A')
}
$$
with $A = \mathcal{A}(U)$, $A' = \mathcal{A}(U')$
is the open immersion of Lemma \ref{lemma-proj-inclusion} above.
\end{enumerate}
\end{lemma}

\begin{proof}
Follows immediately from
Lemmas \ref{lemma-relative-glueing},
\ref{lemma-proj-inclusion}, and
\ref{lemma-transitive-proj}.
\end{proof}

\begin{lemma}
\label{lemma-glue-relative-proj-twists}
In Situation \ref{situation-relative-proj}.
The morphism $\pi : \underline{\text{Proj}}_S(\mathcal{A}) \to S$
of Lemma \ref{lemma-glue-relative-proj} comes with the following
additional structure.
There exists a quasi-coherent $\mathbf{Z}$-graded sheaf
of $\mathcal{O}_{\underline{\text{Proj}}_S(\mathcal{A})}$-algebras
$\bigoplus\nolimits_{n \in \mathbf{Z}}
\mathcal{O}_{\underline{\text{Proj}}_S(\mathcal{A})}(n)$,
and a morphism of graded $\mathcal{O}_S$-algebras
$$
\psi :
\mathcal{A}
\longrightarrow
\bigoplus\nolimits_{n \geq 0}
\pi_*\left(\mathcal{O}_{\underline{\text{Proj}}_S(\mathcal{A})}(n)\right)
$$
uniquely determined by the following property:
For every affine open $U \subset S$ with $A = \mathcal{A}(U)$
there is an isomorphism
$$
\theta_U :
i_U^*\left(
\bigoplus\nolimits_{n \in \mathbf{Z}} \mathcal{O}_{\text{Proj}(A)}(n)
\right)
\longrightarrow
\left(
\bigoplus\nolimits_{n \in \mathbf{Z}}
\mathcal{O}_{\underline{\text{Proj}}_S(\mathcal{A})}(n)
\right)|_{\pi^{-1}(U)}
$$
of $\mathbf{Z}$-graded $\mathcal{O}_{\pi^{-1}(U)}$-algebras
such that
$$
\xymatrix{
A_n
\ar[rr]_\psi
\ar[dr]_-{(\ref{equation-global-sections})}
& &
\Gamma(\pi^{-1}(U),
\mathcal{O}_{\underline{\text{Proj}}_S(\mathcal{A})}(n)) \\
&
\Gamma(\text{Proj}(A),
\mathcal{O}_{\text{Proj}(A)}(n))
\ar[ru]_-{\theta_U}
&
}
$$
is commutative.
\end{lemma}

\begin{proof}
We are going to use Lemma \ref{lemma-relative-glueing-sheaves}
to glue the sheaves of $\mathbf{Z}$-graded algebras
$\bigoplus_{n \in \mathbf{Z}} \mathcal{O}_{\text{Proj}(A)}(n)$
for $A = \mathcal{A}(U)$, $U \subset S$ affine open
over the scheme $\underline{\text{Proj}}_S(\mathcal{A})$.
We have constructed the data necessary for this in
Lemma \ref{lemma-proj-inclusion} and we have checked condition (d) of
Lemma \ref{lemma-relative-glueing-sheaves} in
Lemma \ref{lemma-transitive-proj}. Hence we get the
sheaf of $\mathbf{Z}$-graded
$\mathcal{O}_{\underline{\text{Proj}}_S(\mathcal{A})}$-algebras
$\bigoplus_{n \in \mathbf{Z}}
\mathcal{O}_{\underline{\text{Proj}}_S(\mathcal{A})}(n)$
together with the isomorphisms $\theta_U$ for all
$U \subset S$ affine open and all $n \in \mathbf{Z}$.
For every affine open $U \subset S$ with $A = \mathcal{A}(U)$ we have a map
$A \to \Gamma(\text{Proj}(A),
\bigoplus_{n \geq 0} \mathcal{O}_{\text{Proj}(A)}(n))$.
Hence the map $\psi$ exists by functoriality
of relative glueing, see Remark \ref{remark-relative-glueing-functorial}.
The diagram of the lemma commutes by construction.
This characterizes the sheaf of $\mathbf{Z}$-graded
$\mathcal{O}_{\underline{\text{Proj}}_S(\mathcal{A})}$-algebras
$\bigoplus \mathcal{O}_{\underline{\text{Proj}}_S(\mathcal{A})}(n)$
because the proof of Lemma \ref{lemma-morphism-proj} shows that
having these diagrams commute uniquely determines the maps $\theta_U$.
Some details omitted.
\end{proof}















